\begin{figure}[H]
    \centering
    \begin{tikzpicture}[x=1.0cm, y=1.0cm]

        % (a) the threshold test
        \node[panel, font=\scriptsize\bfseries, anchor=west] at (-1.0, 1.55) {(α) το κατώφλι της ρίζας, $k=4$};

        \node[blkreg, minimum width=0.66cm] (r)  at (1.55,0.85) {$6$};
        \node[blkreg, minimum width=0.66cm] (c1) at (0.85,-0.05) {$8$};
        \node[blkreg, minimum width=0.66cm] (c2) at (2.25,-0.05) {$7$};
        \node[blkreg, minimum width=0.66cm] (c3) at (0.15,-0.95) {$9$};
        \draw[flowmute, -] (r) -- (c1);
        \draw[flowmute, -] (r) -- (c2);
        \draw[flowmute, -] (c1) -- (c3);
        \node[notemute, font=\tiny, anchor=east] at (1.18,0.85) {ρίζα $=$ κατώφλι};
        \node[notemute, font=\tiny, anchor=north, align=center] at (1.55,-1.42)
            {σωρός ελαχίστου με τα\\ \tl{k} καλύτερα (ως τώρα)};

        \node[lanedrop] (x1) at (4.30,0.85) {$3$};
        \node[note, font=\tiny, anchor=west] at (4.68,0.85) {$3 \le 6$};
        \draw[flowmute, dashed, dash pattern=on 2pt off 1.6pt] (x1) -- (r);
        \node[note, font=\scriptsize, anchor=west, text=figregD, align=left] at (6.30,0.85)
            {απόρριψη:\\ \emph{μία} σύγκριση};

        \node[lanereg] (x2) at (4.30,-0.95) {$11$};
        \node[note, font=\tiny, anchor=west] at (4.68,-0.95) {$11 > 6$};
        \draw[flow] (x2) -- (r);
        \node[note, font=\scriptsize, anchor=west, text=figdrmD, align=left] at (6.30,-0.95)
            {αντικατάσταση ρίζας\\ και κατολίσθηση: $O(\log k)$};

        % (b) map-reduce
        \begin{scope}[yshift=-3.55cm]
            \node[panel, font=\scriptsize\bfseries, anchor=west] at (-1.0, 1.10) {(β) το σχήμα \tl{map-reduce}};

            \foreach \i in {0,1,2,3} {
                \pgfmathsetmacro{\xc}{1.05+\i*1.85}
                \node[blk, minimum width=1.55cm, minimum height=0.40cm] (ch\i) at (\xc,0.52)
                    {\tiny τμήμα $\i$};
                \node[blkreg, minimum width=0.85cm, minimum height=0.38cm] (hp\i) at (\xc,-0.42)
                    {\tiny σωρός \tl{k}};
                \draw[flow] (ch\i.south) -- (hp\i.north);
            }
            \node[rowlab, font=\tiny, anchor=west] at (-0.45, 0.52) {\tl{map}};

            \node[blkreg, minimum width=1.55cm, minimum height=0.40cm] (mg) at (3.83,-1.45)
                {\tiny τελικά \tl{k}};
            \foreach \i in {0,1,2,3} { \draw[flowmute] (hp\i.south) -- (mg.north); }
            \node[rowlab, font=\tiny, anchor=west] at (-0.45, -1.45) {\tl{reduce}};

            \node[note, font=\scriptsize, anchor=west, align=left] at (7.70,-0.42)
                {καμία επικοινωνία\\ στη φάση \tl{map}};
            \node[note, font=\scriptsize, anchor=west, align=left, text=figdrmD] at (7.70,-1.45)
                {κάθε σωρός φιλτράρει μόνο\\ αφού δει \tl{k} στοιχεία};
        \end{scope}

    \end{tikzpicture}
    \caption[Επιλογή με σωρούς και σχήμα \tl{map-reduce}]{Η επιλογή με σωρούς και η παραλληλοποίησή της. Στο (α) η ρίζα του
    σωρού ελαχίστου λειτουργεί ως κατώφλι: το στοιχείο που δεν την υπερβαίνει
    απορρίπτεται με μία σύγκριση, ενώ εκείνο που την υπερβαίνει την αντικαθιστά και
    επιβάλλει κατολίσθηση. Στο (β) η είσοδος διαμερίζεται σε τμήματα, ένας σωρός
    ανά εργάτη και τα μερικά αποτελέσματα συγχωνεύονται σε ένα.}
    \label{fig:bg-heap}
\end{figure}
